\clearpage\thispagestyle{empty}
\pdfbookmark[0]{Quellen zu Mondbeobachtungen (Februar 1822)}{moonfebruary1822-sources}
\begin{center}\textbf{Quellen zu Mondbeobachtungen (Februar 1822)}\end{center}
\noindent Diese KI-Übertragung ins heutige Deutsch folgt der ersten Mitteilung mit der Überschrift \textit{Mondsbeobachtungen} in \href{https://gdz.sub.uni-goettingen.de/id/PPN236007467}{Gauß’ \textit{Werken}, Bd.~VI} (Göttingen, 1874), S.~441. H. C. Schumacher veröffentlichte die Mitteilung zuerst unter dem Titel \textit{Sterne, deren grade Aufsteigung mit der des Mondes zu vergleichen ist} in den \textit{Astronomischen Nachrichten}, Bd.~I, \href{https://api.digitale-sammlungen.de/iiif/image/v2/bsb10049426_00061/full/1800,/0/default.jpg}{Nr.~7}, Februar 1822, Sp.~\href{https://api.digitale-sammlungen.de/iiif/image/v2/bsb10049426_00065/full/1800,/0/default.jpg}{105--106}. Er zeichnete mit S. Das \href{https://api.digitale-sammlungen.de/iiif/image/v2/bsb10049426_00005/full/1800,/0/default.jpg}{Titelblatt} des gebundenen Bandes aus der Bayerischen Staatsbibliothek nennt Altona 1823 und die Hammerich- und Heineking’sche Buchdruckerei.
\par\smallskip\noindent Schumacher nennt Professor Encke als Übermittler der Sterne für die nächsten Lunationen und Gauß als Übermittler der Göttinger Beobachtungen für die vorigen. Im Erstdruck stehen einfache und doppelte Striche, wo die Werkausgabe hochgestellte m und s verwendet. In beiden Drucken bleiben die Differenzen in den beiden Mondzeilen leer. KI-Übertragung und KI-Satz.
\clearpage
